In multi-agent satellite--ground scheduling, a coordinated adjustment is valuable
only when its immediate change and recovered contacts form a compatible
schedule. Summing individual replacement values can overestimate improvement,
while solving every candidate can exhaust the decision budget. We formulate
learning the additional joint recovery beyond a shared, executed greedy
response. An action-conditioned graph encoder represents eligible contacts,
conflict factors and actual warm membership. A request-aware head combines
fractional capacity features with warm and budget context to predict finite
recovery between a feasible lower value and a clique-partition upper bound;
the known immediate gain is added separately. Training compares alternative
requests executed from cloned states, using signed values, improvement rankings
and feasible membership supervision. At deployment, predictions prioritize
finite executor calls while an independent kernel checks complete schedules
and preserves the incumbent. Quality measures feasible improvement actually
received before the total deadline, charging preparation, inference, execution
and final checks.
Controlled experiments identify substantial conditional recovery opportunity,
but distinguish it from realized policy quality: longer allowances permit
additional recovery, while inexpensive summary and greedy controls remain
competitive and strict short deadlines expose the cost of allocation.
Profiles of seven published native configurations and three released neural
recipes on 40 public graphs provide a complementary solver comparison.
The results clarify when joint recovery is an informative learning target
and why structural prediction alone does not guarantee better timed schedules.
